\documentclass{article}%
\usepackage[T1]{fontenc}%
\usepackage[utf8]{inputenc}%
\usepackage{lmodern}%
\usepackage{textcomp}%
\usepackage{lastpage}%
%
\usepackage[russian]{babel}%
\usepackage{fontenc}%
\usepackage{graphicx}%
\usepackage{float}%
\usepackage{longtable}%
\usepackage{array}%
%
\begin{document}%
\normalsize%
\begin{longtable}{|>{\centering\arraybackslash}p{17.176750000000002cm}|}%
\hline%
ФГБОУ ВО ЮГОРСКИЙ ГОСУДАРСТВЕННЫЙ УНИВЕРСИТЕТ% Шрифт-Times New Roman Размер шрифта-14.0 Межстрочный интервал-1.5 Правило-ONE_POINT_FIVE (1) %
 \\ \hline \endfirsthead%
\hline%
ИНЖЕНЕРНАЯ ШКОЛА ЦИФРОВЫХ ТЕХНОЛОГИЙ% Шрифт-Times New Roman Размер шрифта-14.0 Межстрочный интервал-1.5 Правило-ONE_POINT_FIVE (1) %
 \\ \hline%
\end{longtable}%
\begin{center}%
ОТЧЕТ% Шрифт-Times New Roman Размер шрифта-14.0 Межстрочный интервал-1.5 Правило-ONE_POINT_FIVE (1) %
%
\end{center}%
\begin{center}%
по лабораторной работе № 2:% Шрифт-Times New Roman Размер шрифта-14.0 Межстрочный интервал-1.5 Правило-ONE_POINT_FIVE (1) %
%
\end{center}%
\begin{center}%
«Метод наискорейшего градиентного спуска»% Шрифт-Times New Roman Размер шрифта-14.0 Межстрочный интервал-1.5 Правило-ONE_POINT_FIVE (1) %
%
\end{center}%
\begin{center}%
По дисциплине:% Шрифт-Times New Roman Размер шрифта-14.0 Межстрочный интервал-1.5 Правило-ONE_POINT_FIVE (1) %
%
\end{center}%
\begin{center}%
«Машинное обучение»% Шрифт-Times New Roman Размер шрифта-14.0 Межстрочный интервал-1.5 Правило-ONE_POINT_FIVE (1) %
%
\end{center}%
\begin{flushleft}%
Выполнил: студент группы 1521б% Шрифт-Times New Roman Размер шрифта-14.0 Межстрочный интервал-1.5 Правило-ONE_POINT_FIVE (1) %
%
\end{flushleft}%
\begin{flushleft}%
Евдакимов Егор Александрович% Шрифт-Times New Roman Размер шрифта-14.0 Межстрочный интервал-1.5 Правило-ONE_POINT_FIVE (1) %
%
\end{flushleft}%
\begin{flushleft}%
Проверил: старший преподаватель% Шрифт-Times New Roman Размер шрифта-14.0 Межстрочный интервал-1.5 Правило-ONE_POINT_FIVE (1) %
%
\end{flushleft}%
\begin{flushleft}%
Шицелов Анатолий Вячеславович% Шрифт-Times New Roman Размер шрифта-14.0 Межстрочный интервал-1.5 Правило-ONE_POINT_FIVE (1) %
%
\end{flushleft}%
\begin{center}%
г. Ханты-Мансийск, 2025 г.% Шрифт-Times New Roman Размер шрифта-14.0 Межстрочный интервал-1.5 Правило-ONE_POINT_FIVE (1) %
%
\end{center}%
\begin{center}%
\textbf{Содержание}% Шрифт-Times New Roman Размер шрифта-14.0 Межстрочный интервал-1.5 Правило-ONE_POINT_FIVE (1) %
%
\end{center}%
\begin{flushleft}%
% Шрифт-Times New Roman Размер шрифта-14.0 Межстрочный интервал-None Правило-None %
%
\end{flushleft}%
\section{Описание работы}%
\begin{flushleft}%
	% Шрифт-Times New Roman Размер шрифта-14.0 Межстрочный интервал-1.0 Правило-SINGLE (0) %
\textbf{Цель работы}% Шрифт-None Размер шрифта-12.0 Межстрочный интервал-1.0 Правило-SINGLE (0) %
: реализовать метод наискорейшего градиентного спуска для нахождения точки глобального минимума функции Бута.% Шрифт-None Размер шрифта-12.0 Межстрочный интервал-1.0 Правило-SINGLE (0) %
%
\end{flushleft}%
\begin{center}%
\textbf{Задачи}% Шрифт-Times New Roman Размер шрифта-14.0 Межстрочный интервал-1.5 Правило-ONE_POINT_FIVE (1) %
%
\end{center}%
\begin{enumerate}%
\item Реализовать метода наискорейшего градиентного спуска% Шрифт-Times New Roman Размер шрифта-14.0 Межстрочный интервал-1.5 Правило-ONE_POINT_FIVE (1) %
%
\item Отобразить график изменения значения функции по шагам% Шрифт-Times New Roman Размер шрифта-14.0 Межстрочный интервал-1.5 Правило-ONE_POINT_FIVE (1) %
%
\item Отобразить 3-х мерный график функции от X и Y% Шрифт-Times New Roman Размер шрифта-14.0 Межстрочный интервал-1.5 Правило-ONE_POINT_FIVE (1) %
%
\item Отобразить движение алгоритма на 3-х мерном графике% Шрифт-Times New Roman Размер шрифта-14.0 Межстрочный интервал-1.5 Правило-ONE_POINT_FIVE (1) %
%
\end{enumerate}%
\subsection{Ход работы}%
\begin{flushleft}%
Формулы, которые понадобятся для выполнения работы:% Шрифт-Times New Roman Размер шрифта-14.0 Межстрочный интервал-1.5 Правило-ONE_POINT_FIVE (1) %
%
\end{flushleft}%
\begin{enumerate}%
\item Функция Бута % Шрифт-Times New Roman Размер шрифта-14.0 Межстрочный интервал-1.0 Правило-SINGLE (0) %
%
\item %
\item Частная производная по x% Шрифт-Times New Roman Размер шрифта-14.0 Межстрочный интервал-None Правило-None %
%
\item %
\item Частная производная по y% Шрифт-Times New Roman Размер шрифта-14.0 Межстрочный интервал-None Правило-None %
%
\item %
\begin{flushleft}%
В ходе реализации лабораторной работы были импортированы следующие модули Python (Листинг 1):% Шрифт-Times New Roman Размер шрифта-14.0 Межстрочный интервал-1.5 Правило-ONE_POINT_FIVE (1) %
%
\end{flushleft}%
\begin{flushleft}%
Листинг 1 – Импорт модулей% Шрифт-Times New Roman Размер шрифта-14.0 Межстрочный интервал-1.5 Правило-ONE_POINT_FIVE (1) %
%
\end{flushleft}%
\begin{flushleft}%
import math% Шрифт-Courier New Размер шрифта-10.0 Межстрочный интервал-1.0 Правило-SINGLE (0) %
%
\end{flushleft}%
\begin{flushleft}%
import numpy as np% Шрифт-Courier New Размер шрифта-10.0 Межстрочный интервал-1.0 Правило-SINGLE (0) %
%
\end{flushleft}%
\begin{flushleft}%
import matplotlib.pyplot as plt% Шрифт-Courier New Размер шрифта-10.0 Межстрочный интервал-1.0 Правило-SINGLE (0) %
%
\end{flushleft}%
\begin{flushleft}%
from mpl\textbackslash{}_toolkits.mplot3d import Axes3D% Шрифт-Courier New Размер шрифта-10.0 Межстрочный интервал-1.0 Правило-SINGLE (0) %
%
\end{flushleft}%
\begin{flushleft}%
Реализуем функцию Бута (Листинг 2):% Шрифт-Times New Roman Размер шрифта-14.0 Межстрочный интервал-1.5 Правило-ONE_POINT_FIVE (1) %
%
\end{flushleft}%
\begin{flushleft}%
Листинг 2 – Функция Бута и ее частные производные% Шрифт-Times New Roman Размер шрифта-14.0 Межстрочный интервал-1.5 Правило-ONE_POINT_FIVE (1) %
%
\end{flushleft}%
\begin{flushleft}%
def f(x, y):% Шрифт-Courier New Размер шрифта-10.0 Межстрочный интервал-1.0 Правило-SINGLE (0) %
%
\end{flushleft}%
\begin{flushleft}%
    return (x + 2 * y - 7) ** 2 + (2 * x + y - 5) ** 2% Шрифт-Courier New Размер шрифта-10.0 Межстрочный интервал-1.0 Правило-SINGLE (0) %
%
\end{flushleft}%
\begin{flushleft}%
def dx(x, y):% Шрифт-Courier New Размер шрифта-10.0 Межстрочный интервал-1.0 Правило-SINGLE (0) %
%
\end{flushleft}%
\begin{flushleft}%
    return 10 * x + 8 * y - 34% Шрифт-Courier New Размер шрифта-10.0 Межстрочный интервал-1.0 Правило-SINGLE (0) %
%
\end{flushleft}%
\begin{flushleft}%
def dy(x, y):% Шрифт-Courier New Размер шрифта-10.0 Межстрочный интервал-1.0 Правило-SINGLE (0) %
%
\end{flushleft}%
\begin{flushleft}%
    return 8 * x + 10 * y – 38% Шрифт-Courier New Размер шрифта-10.0 Межстрочный интервал-1.0 Правило-SINGLE (0) %
%
\end{flushleft}%
\begin{flushleft}%
В этом листинге определена целевая функция Бута, которая используется для тестирования оптимизационных алгоритмов.
Функция имеет глобальный минимум в точке (1,3).
Также вычисляются ее частные производные, необходимые для градиентного спуска.% Шрифт-Times New Roman Размер шрифта-14.0 Межстрочный интервал-1.5 Правило-ONE_POINT_FIVE (1) %
%
\end{flushleft}%
\begin{flushleft}%
	Реализуем функцию золотого сечения для дальнейшего динамического вычисления шага обучения (Листинг 3):% Шрифт-Times New Roman Размер шрифта-14.0 Межстрочный интервал-1.5 Правило-ONE_POINT_FIVE (1) %
%
\end{flushleft}%
\begin{flushleft}%
	Листинг 3 – Функция золотого сечения% Шрифт-Times New Roman Размер шрифта-14.0 Межстрочный интервал-1.5 Правило-ONE_POINT_FIVE (1) %
%
\end{flushleft}%
\begin{flushleft}%
def golden\textbackslash{}_ratio(f, x, y, grad\textbackslash{}_x, grad\textbackslash{}_y):% Шрифт-Courier New Размер шрифта-10.0 Межстрочный интервал-1.0 Правило-SINGLE (0) %
%
\end{flushleft}%
\begin{flushleft}%
    iter=100% Шрифт-Courier New Размер шрифта-10.0 Межстрочный интервал-1.0 Правило-SINGLE (0) %
%
\end{flushleft}%
\begin{flushleft}%
    tol=1e-5% Шрифт-Courier New Размер шрифта-10.0 Межстрочный интервал-1.0 Правило-SINGLE (0) %
%
\end{flushleft}%
\begin{flushleft}%
    a = 0% Шрифт-Courier New Размер шрифта-10.0 Межстрочный интервал-1.0 Правило-SINGLE (0) %
%
\end{flushleft}%
\begin{flushleft}%
    b = 1% Шрифт-Courier New Размер шрифта-10.0 Межстрочный интервал-1.0 Правило-SINGLE (0) %
%
\end{flushleft}%
\begin{flushleft}%
    fi = (math.sqrt(5) + 1) / 2% Шрифт-Courier New Размер шрифта-10.0 Межстрочный интервал-1.0 Правило-SINGLE (0) %
%
\end{flushleft}%
\begin{flushleft}%
    c = b - (b - a) / fi% Шрифт-Courier New Размер шрифта-10.0 Межстрочный интервал-1.0 Правило-SINGLE (0) %
%
\end{flushleft}%
\begin{flushleft}%
    d = a + (b - a) / fi% Шрифт-Courier New Размер шрифта-10.0 Межстрочный интервал-1.0 Правило-SINGLE (0) %
%
\end{flushleft}%
\begin{flushleft}%
    for \textbackslash{}_ in range(iter):% Шрифт-Courier New Размер шрифта-10.0 Межстрочный интервал-1.0 Правило-SINGLE (0) %
%
\end{flushleft}%
\begin{flushleft}%
        fc = f(x - c * grad\textbackslash{}_x, y - c * grad\textbackslash{}_y)% Шрифт-Courier New Размер шрифта-10.0 Межстрочный интервал-1.0 Правило-SINGLE (0) %
%
\end{flushleft}%
\begin{flushleft}%
        fd = f(x - d * grad\textbackslash{}_x, y - d * grad\textbackslash{}_y)% Шрифт-Courier New Размер шрифта-10.0 Межстрочный интервал-1.0 Правило-SINGLE (0) %
%
\end{flushleft}%
\begin{flushleft}%
        if fc $<$ fd:% Шрифт-Courier New Размер шрифта-10.0 Межстрочный интервал-1.0 Правило-SINGLE (0) %
%
\end{flushleft}%
\begin{flushleft}%
            b = d% Шрифт-Courier New Размер шрифта-10.0 Межстрочный интервал-1.0 Правило-SINGLE (0) %
%
\end{flushleft}%
\begin{flushleft}%
        else:% Шрифт-Courier New Размер шрифта-10.0 Межстрочный интервал-1.0 Правило-SINGLE (0) %
%
\end{flushleft}%
\begin{flushleft}%
            a = c% Шрифт-Courier New Размер шрифта-10.0 Межстрочный интервал-1.0 Правило-SINGLE (0) %
%
\end{flushleft}%
\begin{flushleft}%
        if abs(b - a) $<$ tol:% Шрифт-Courier New Размер шрифта-10.0 Межстрочный интервал-1.0 Правило-SINGLE (0) %
%
\end{flushleft}%
\begin{flushleft}%
            break  % Шрифт-Courier New Размер шрифта-10.0 Межстрочный интервал-1.0 Правило-SINGLE (0) %
%
\end{flushleft}%
\begin{flushleft}%
        c = b - (b - a) / fi% Шрифт-Courier New Размер шрифта-10.0 Межстрочный интервал-1.0 Правило-SINGLE (0) %
%
\end{flushleft}%
\begin{flushleft}%
        d = a + (b - a) / fi% Шрифт-Courier New Размер шрифта-10.0 Межстрочный интервал-1.0 Правило-SINGLE (0) %
%
\end{flushleft}%
\begin{flushleft}%
    return (a + b) / 2% Шрифт-Courier New Размер шрифта-10.0 Межстрочный интервал-1.0 Правило-SINGLE (0) %
%
\end{flushleft}%
\begin{flushleft}%
Метод % Шрифт-Times New Roman Размер шрифта-14.0 Межстрочный интервал-1.5 Правило-ONE_POINT_FIVE (1) %
золотого сечения% Шрифт-Times New Roman Размер шрифта-14.0 Межстрочный интервал-1.5 Правило-ONE_POINT_FIVE (1) %
, он используется для нахождения оптимального шага градиентного спуска.
Метод итеративно сужает диапазон возможных значений шага, пока не достигнет заданной точности (tol=1e-5).
В ходе работы алгоритм обновляет границы интервала [a, b] и вычисляет значения функции в точках % Шрифт-Times New Roman Размер шрифта-14.0 Межстрочный интервал-1.5 Правило-ONE_POINT_FIVE (1) %
c% Шрифт-Times New Roman Размер шрифта-14.0 Межстрочный интервал-1.5 Правило-ONE_POINT_FIVE (1) %
 % Шрифт-Times New Roman Размер шрифта-14.0 Межстрочный интервал-1.5 Правило-ONE_POINT_FIVE (1) %
и % Шрифт-Times New Roman Размер шрифта-14.0 Межстрочный интервал-1.5 Правило-ONE_POINT_FIVE (1) %
d% Шрифт-Times New Roman Размер шрифта-14.0 Межстрочный интервал-1.5 Правило-ONE_POINT_FIVE (1) %
, пока не найдет наилучший шаг.% Шрифт-Times New Roman Размер шрифта-14.0 Межстрочный интервал-1.5 Правило-ONE_POINT_FIVE (1) %
%
\end{flushleft}%
\begin{flushleft}%
	% Шрифт-Times New Roman Размер шрифта-14.0 Межстрочный интервал-1.5 Правило-ONE_POINT_FIVE (1) %
%
\end{flushleft}%
\begin{flushleft}%
% Шрифт-Times New Roman Размер шрифта-14.0 Межстрочный интервал-None Правило-None %
%
\end{flushleft}%
\begin{flushleft}%
Реализуем функцию градиентного спуска (Листинг 4):% Шрифт-Times New Roman Размер шрифта-14.0 Межстрочный интервал-1.5 Правило-ONE_POINT_FIVE (1) %
%
\end{flushleft}%
\begin{flushleft}%
	Листинг 4 – Функция градиентного спуска% Шрифт-Times New Roman Размер шрифта-14.0 Межстрочный интервал-1.5 Правило-ONE_POINT_FIVE (1) %
%
\end{flushleft}%
\begin{flushleft}%
def gradient\textbackslash{}_descent():% Шрифт-Courier New Размер шрифта-10.0 Межстрочный интервал-1.0 Правило-SINGLE (0) %
%
\end{flushleft}%
\begin{flushleft}%
    num\textbackslash{}_iter = 100% Шрифт-Courier New Размер шрифта-10.0 Межстрочный интервал-1.0 Правило-SINGLE (0) %
%
\end{flushleft}%
\begin{flushleft}%
    x = np.random.uniform(-10, 10)% Шрифт-Courier New Размер шрифта-10.0 Межстрочный интервал-1.0 Правило-SINGLE (0) %
%
\end{flushleft}%
\begin{flushleft}%
    y = np.random.uniform(-10, 10)% Шрифт-Courier New Размер шрифта-10.0 Межстрочный интервал-1.0 Правило-SINGLE (0) %
%
\end{flushleft}%
\begin{flushleft}%
    history = [(x, y, f(x, y))]% Шрифт-Courier New Размер шрифта-10.0 Межстрочный интервал-1.0 Правило-SINGLE (0) %
%
\end{flushleft}%
\begin{flushleft}%
       for \textbackslash{}_ in range(num\textbackslash{}_iter):% Шрифт-Courier New Размер шрифта-10.0 Межстрочный интервал-1.0 Правило-SINGLE (0) %
%
\end{flushleft}%
\begin{flushleft}%
        grad\textbackslash{}_x = dx(x, y)% Шрифт-Courier New Размер шрифта-10.0 Межстрочный интервал-1.0 Правило-SINGLE (0) %
%
\end{flushleft}%
\begin{flushleft}%
        grad\textbackslash{}_y = dy(x, y)% Шрифт-Courier New Размер шрифта-10.0 Межстрочный интервал-1.0 Правило-SINGLE (0) %
%
\end{flushleft}%
\begin{flushleft}%
        alpha = golden\textbackslash{}_ratio(f, x, y, grad\textbackslash{}_x, grad\textbackslash{}_y)% Шрифт-Courier New Размер шрифта-10.0 Межстрочный интервал-1.0 Правило-SINGLE (0) %
%
\end{flushleft}%
\begin{flushleft}%
        x\textbackslash{}_2 = x - alpha * grad\textbackslash{}_x% Шрифт-Courier New Размер шрифта-10.0 Межстрочный интервал-1.0 Правило-SINGLE (0) %
%
\end{flushleft}%
\begin{flushleft}%
        y\textbackslash{}_2 = y - alpha * grad\textbackslash{}_y % Шрифт-Courier New Размер шрифта-10.0 Межстрочный интервал-1.0 Правило-SINGLE (0) %
%
\end{flushleft}%
\begin{flushleft}%
        check\textbackslash{}_x = x\textbackslash{}_2 - x% Шрифт-Courier New Размер шрифта-10.0 Межстрочный интервал-1.0 Правило-SINGLE (0) %
%
\end{flushleft}%
\begin{flushleft}%
        check\textbackslash{}_y = y\textbackslash{}_2 - y% Шрифт-Courier New Размер шрифта-10.0 Межстрочный интервал-1.0 Правило-SINGLE (0) %
%
\end{flushleft}%
\begin{flushleft}%
        check\textbackslash{}_xy = np.sqrt(check\textbackslash{}_x ** 2 + check\textbackslash{}_y ** 2)% Шрифт-Courier New Размер шрифта-10.0 Межстрочный интервал-1.0 Правило-SINGLE (0) %
%
\end{flushleft}%
\begin{flushleft}%
        if check\textbackslash{}_xy $<$ 0.00001: break% Шрифт-Courier New Размер шрифта-10.0 Межстрочный интервал-1.0 Правило-SINGLE (0) %
%
\end{flushleft}%
\begin{flushleft}%
        x, y = x\textbackslash{}_2, y\textbackslash{}_2% Шрифт-Courier New Размер шрифта-10.0 Межстрочный интервал-1.0 Правило-SINGLE (0) %
%
\end{flushleft}%
\begin{flushleft}%
        history.append((x, y, f(x, y)))% Шрифт-Courier New Размер шрифта-10.0 Межстрочный интервал-1.0 Правило-SINGLE (0) %
%
\end{flushleft}%
\begin{flushleft}%
    return x, y, f(x, y), history% Шрифт-Courier New Размер шрифта-10.0 Межстрочный интервал-1.0 Правило-SINGLE (0) %
%
\end{flushleft}%
\begin{flushleft}%
Алгоритм % Шрифт-Times New Roman Размер шрифта-14.0 Межстрочный интервал-1.5 Правило-ONE_POINT_FIVE (1) %
наискорейшего градиентного спуска% Шрифт-Times New Roman Размер шрифта-14.0 Межстрочный интервал-1.5 Правило-ONE_POINT_FIVE (1) %
, который находит минимум функции Бута. Он генерирует случайные стартовые значения, вычисляет градиент в данной точке, ищет оптимальный шаг с помощью функции золотого сечения, вычисляет следующую точку, от которой будет считаться следующий градиент и сохраняет все точки в массиве для дальнейшей визуализации спуска, также добавим остановку метода если точность найденного минимума меньше 0.00001.% Шрифт-Times New Roman Размер шрифта-14.0 Межстрочный интервал-1.5 Правило-ONE_POINT_FIVE (1) %
%
\end{flushleft}%
\begin{flushleft}%
Вот, что мы получили при запуске кода (рисунок 1, рисунок 2, рисунок 3):% Шрифт-Times New Roman Размер шрифта-14.0 Межстрочный интервал-1.5 Правило-ONE_POINT_FIVE (1) %
%
\end{flushleft}%
\begin{flushleft}%
Вывод в консоль:% Шрифт-Times New Roman Размер шрифта-14.0 Межстрочный интервал-1.5 Правило-ONE_POINT_FIVE (1) %
%
\end{flushleft}%
\begin{flushleft}%
Точка минимума x: 0.9999999999999996 y: 3.0000000000000004% Шрифт-Times New Roman Размер шрифта-14.0 Межстрочный интервал-1.5 Правило-ONE_POINT_FIVE (1) %
%
\end{flushleft}%
\begin{flushleft}%
Глобальный минимум функции Бута (1, 3), это практически совпадает с % Шрифт-Times New Roman Размер шрифта-14.0 Межстрочный интервал-1.5 Правило-ONE_POINT_FIVE (1) %
найденным значением% Шрифт-Times New Roman Размер шрифта-14.0 Межстрочный интервал-1.5 Правило-ONE_POINT_FIVE (1) %
\textbf{,}% Шрифт-Times New Roman Размер шрифта-14.0 Межстрочный интервал-1.5 Правило-ONE_POINT_FIVE (1) %
 что подтверждает корректность реализации метода.% Шрифт-Times New Roman Размер шрифта-14.0 Межстрочный интервал-1.5 Правило-ONE_POINT_FIVE (1) %
%
\end{flushleft}%
\end{enumerate}%
\begin{figure}[H]%
\centering%
\includegraphics[width=0.8\textwidth]{images/img\_6.png}%
\end{figure}%
\begin{center}%
Рисунок 1 – 3-х мерный график функции Билла% Шрифт-Times New Roman Размер шрифта-14.0 Межстрочный интервал-1.5 Правило-ONE_POINT_FIVE (1) %
%
\end{center}%
\begin{flushleft}%
На данном графике представлена поверхность функции Бута в трехмерном пространстве. X и Y – входные параметры функции; Z – значение функции в соответствующей точке.
По данному графику можно наблюдать форму целевой функции, наличие минимумов и характер ее изменения.% Шрифт-Times New Roman Размер шрифта-14.0 Межстрочный интервал-1.5 Правило-ONE_POINT_FIVE (1) %
%
\end{flushleft}%
\begin{flushleft}%
	% Шрифт-Times New Roman Размер шрифта-14.0 Межстрочный интервал-None Правило-None %
%
\end{flushleft}%
\begin{flushleft}%
% Шрифт-Times New Roman Размер шрифта-14.0 Межстрочный интервал-None Правило-None %
%
\end{flushleft}%
\begin{figure}[H]%
\centering%
\includegraphics[width=0.8\textwidth]{images/img\_7.png}%
\end{figure}%
\begin{center}%
Рисунок 2 – движение алгоритма на 3-х мерном графике% Шрифт-Times New Roman Размер шрифта-14.0 Межстрочный интервал-None Правило-None %
%
\end{center}%
\begin{flushleft}%
Этот рисунок демонстрирует % Шрифт-Times New Roman Размер шрифта-14.0 Межстрочный интервал-1.5 Правило-ONE_POINT_FIVE (1) %
траекторию движения алгоритма% Шрифт-Times New Roman Размер шрифта-14.0 Межстрочный интервал-1.5 Правило-ONE_POINT_FIVE (1) %
 градиентного спуска по поверхности функции.
Красная линия отображает путь, который проделывает алгоритм от случайной начальной точки к глобальному минимуму.
По визуализации видно, что алгоритм движется вниз по поверхности, постепенно снижая значение функции.% Шрифт-Times New Roman Размер шрифта-14.0 Межстрочный интервал-1.5 Правило-ONE_POINT_FIVE (1) %
%
\end{flushleft}%
\begin{flushleft}%
	% Шрифт-Times New Roman Размер шрифта-14.0 Межстрочный интервал-1.5 Правило-ONE_POINT_FIVE (1) %
%
\end{flushleft}%
\begin{flushleft}%
% Шрифт-Times New Roman Размер шрифта-14.0 Межстрочный интервал-None Правило-None %
%
\end{flushleft}%
\begin{figure}[H]%
\centering%
\includegraphics[width=0.8\textwidth]{images/img\_8.png}%
\end{figure}%
\begin{center}%
Рисунок 3 – график изменения значения функции по шагам% Шрифт-Times New Roman Размер шрифта-14.0 Межстрочный интервал-1.5 Правило-ONE_POINT_FIVE (1) %
%
\end{center}%
\begin{flushleft}%
	Этот график показывает, как % Шрифт-None Размер шрифта-14.0 Межстрочный интервал-1.5 Правило-ONE_POINT_FIVE (1) %
значение целевой функции% Шрифт-None Размер шрифта-14.0 Межстрочный интервал-1.5 Правило-ONE_POINT_FIVE (1) %
 % Шрифт-None Размер шрифта-14.0 Межстрочный интервал-1.5 Правило-ONE_POINT_FIVE (1) %
уменьшается% Шрифт-None Размер шрифта-14.0 Межстрочный интервал-1.5 Правило-ONE_POINT_FIVE (1) %
 с каждой итерацией алгоритма.
На оси % Шрифт-None Размер шрифта-14.0 Межстрочный интервал-1.5 Правило-ONE_POINT_FIVE (1) %
X% Шрифт-None Размер шрифта-14.0 Межстрочный интервал-1.5 Правило-ONE_POINT_FIVE (1) %
 – номер итерации, на оси % Шрифт-None Размер шрифта-14.0 Межстрочный интервал-1.5 Правило-ONE_POINT_FIVE (1) %
Y% Шрифт-None Размер шрифта-14.0 Межстрочный интервал-1.5 Правило-ONE_POINT_FIVE (1) %
 – значение функции Бута.
По графику видно, что алгоритм быстро снижает значение функции на первых итерациях, а затем приближается к минимальному значению, показывая % Шрифт-None Размер шрифта-14.0 Межстрочный интервал-1.5 Правило-ONE_POINT_FIVE (1) %
сходимость метода% Шрифт-None Размер шрифта-14.0 Межстрочный интервал-1.5 Правило-ONE_POINT_FIVE (1) %
\textbf{.}% Шрифт-None Размер шрифта-14.0 Межстрочный интервал-1.5 Правило-ONE_POINT_FIVE (1) %
%
\end{flushleft}%
\begin{flushleft}%
Так же вычислим mape(средняя абсолютная ошибка)(Листинг 5, Рисунок 4)% Шрифт-None Размер шрифта-14.0 Межстрочный интервал-1.5 Правило-ONE_POINT_FIVE (1) %
%
\end{flushleft}%
\begin{flushleft}%
Листинг 5 – Функция mape% Шрифт-None Размер шрифта-14.0 Межстрочный интервал-1.5 Правило-ONE_POINT_FIVE (1) %
%
\end{flushleft}%
\begin{flushleft}%
def mape(history, optimal\textbackslash{}_point):% Шрифт-Courier New Размер шрифта-10.0 Межстрочный интервал-None Правило-None %
%
\end{flushleft}%
\begin{flushleft}%
    final\textbackslash{}_mape = 0% Шрифт-Courier New Размер шрифта-10.0 Межстрочный интервал-None Правило-None %
%
\end{flushleft}%
\begin{flushleft}%
    for point in history:% Шрифт-Courier New Размер шрифта-10.0 Межстрочный интервал-None Правило-None %
%
\end{flushleft}%
\begin{flushleft}%
        m\textbackslash{}_x = abs((optimal\textbackslash{}_point[0] - point[0]) / optimal\textbackslash{}_point[0]) * 100% Шрифт-Courier New Размер шрифта-10.0 Межстрочный интервал-None Правило-None %
%
\end{flushleft}%
\begin{flushleft}%
        m\textbackslash{}_y = abs((optimal\textbackslash{}_point[1] - point[1]) / optimal\textbackslash{}_point[1]) * 100% Шрифт-Courier New Размер шрифта-10.0 Межстрочный интервал-None Правило-None %
%
\end{flushleft}%
\begin{flushleft}%
        mape = (m\textbackslash{}_x + m\textbackslash{}_y) / 2% Шрифт-Courier New Размер шрифта-10.0 Межстрочный интервал-None Правило-None %
%
\end{flushleft}%
\begin{flushleft}%
        final\textbackslash{}_mape += mape% Шрифт-Courier New Размер шрифта-10.0 Межстрочный интервал-None Правило-None %
%
\end{flushleft}%
\begin{flushleft}%
    final\textbackslash{}_mape /= len(history)% Шрифт-Courier New Размер шрифта-10.0 Межстрочный интервал-None Правило-None %
%
\end{flushleft}%
\begin{flushleft}%
    return final\textbackslash{}_mape% Шрифт-Courier New Размер шрифта-10.0 Межстрочный интервал-None Правило-None %
%
\end{flushleft}%
\begin{flushleft}%
final\textbackslash{}_mape = mape(history, optimal\textbackslash{}_point)% Шрифт-Courier New Размер шрифта-10.0 Межстрочный интервал-None Правило-None %
%
\end{flushleft}%
\begin{flushleft}%
print(f'Средняя абсолютная ошибка: \textbackslash{}{final\textbackslash{}_mape:.3f\textbackslash{}} \textbackslash{}%')% Шрифт-Courier New Размер шрифта-10.0 Межстрочный интервал-None Правило-None %
%
\end{flushleft}%
\begin{flushleft}%
В данном листинге реализована функция mape, предназначенная для оценки точности работы алгоритма градиентного спуска. MAPE (Mean Absolute Percentage Error) вычисляется как среднее значение процентной ошибки по осям X и Y на каждой итерации, по отношению к оптимальной точке минимума (в данном случае — (1, 3)). Итоговое значение MAPE позволяет количественно оценить, насколько эффективно алгоритм сходится к оптимальному решению, и может использоваться для сравнения с другими методами оптимизации. Чем ниже значение MAPE, тем ближе траектория спуска к реальному минимуму функции.% Шрифт-None Размер шрифта-14.0 Межстрочный интервал-1.5 Правило-ONE_POINT_FIVE (1) %
%
\end{flushleft}%
\begin{figure}[H]%
\centering%
\includegraphics[width=0.8\textwidth]{images/img\_9.png}%
\end{figure}%
\begin{center}%
Рисунок 4 – Средняя абсолютная ошибка% Шрифт-Times New Roman Размер шрифта-14.0 Межстрочный интервал-1.5 Правило-ONE_POINT_FIVE (1) %
%
\end{center}%
\begin{flushleft}%
На изображении представлен текстовый вывод, демонстрирующий итоговое значение средней абсолютной процентной ошибки (MAPE), равное 4.585 \textbackslash{}%. Этот показатель отражает среднюю относительную погрешность приближенного решения, полученного методом градиентного спуска. Низкое значение ошибки подтверждает, что алгоритм с динамически изменяющимся шагом эффективно сходится к точке минимума функции.% Шрифт-Times New Roman Размер шрифта-14.0 Межстрочный интервал-1.5 Правило-ONE_POINT_FIVE (1) %
% Шрифт-Times New Roman Размер шрифта-14.0 Межстрочный интервал-1.5 Правило-ONE_POINT_FIVE (1) %
%
\end{flushleft}%
\section{Заключение}%
\begin{flushleft}%
В ходе выполнения лабораторной работы были успешно решены все поставленные задачи.% Шрифт-Times New Roman Размер шрифта-14.0 Межстрочный интервал-1.5 Правило-ONE_POINT_FIVE (1) %
%
\end{flushleft}%
\begin{flushleft}%
Был разработан и реализован метод наискорейшего градиентного спуска для поиска глобального минимума функции Бута. В процессе работы алгоритм динамически вычислял оптимальный шаг с помощью метода золотого сечения для ускорения сходимости.% Шрифт-Times New Roman Размер шрифта-14.0 Межстрочный интервал-1.5 Правило-ONE_POINT_FIVE (1) %
%
\end{flushleft}%
\begin{flushleft}%
Для визуализации результатов были построены следующие графики:3D-график функции Бута, траектория движения алгоритма в трехмерном пространстве, график изменения значения функции по итерациям.% Шрифт-Times New Roman Размер шрифта-14.0 Межстрочный интервал-1.5 Правило-ONE_POINT_FIVE (1) %
%
\end{flushleft}%
\begin{flushleft}%
Результаты работы алгоритма показали, что он успешно находит глобальный минимум функции в точке (1,3) за наименьшее количество итераций, что подтверждает его эффективность.% Шрифт-Times New Roman Размер шрифта-14.0 Межстрочный интервал-1.5 Правило-ONE_POINT_FIVE (1) %
%
\end{flushleft}%
\end{document}